\begin{abstract}
A drowsiness detector in a vehicle must work on a new driver without a calibration session. Online Euclidean alignment offers this, and it is now being used in electroencephalography (EEG) test-time adaptation, but its behaviour during a real drive is poorly understood. Using all 62 sessions from 27 drivers of a public lane-keeping dataset under leave-one-subject-out testing, we examine when it helps. Short references are near-singular in every session and collapse accuracy to 52.55\%; shrinkage restores it to 74.79\%. Blending the running covariance with a population prior raises balanced accuracy from 73.59\% to 77.37\% ($p = 0.039$), matching whole-session alignment, although the result depends on a few drivers and AUC does not improve. Exploratory analysis shows the gain varies with the drowsy share of windows seen so far: about one point at low share, 7.70 and 12.47 points in the middle bins, and about one point above 75\%. This pattern is entangled with time on task. A frozen early reference loses 6.40 points late in the drive. Updating batch-normalisation statistics during the drive harms this model, lowering accuracy to 53.84\% at a published setting. Alignment also increases reliance on gamma and theta activity and weakens attribution tests. These results come from one model and one dataset, but they suggest that the mix of states in the reference, not only its length, should guide online EEG adaptation.
\end{abstract}

\begin{IEEEkeywords}
Drowsiness detection, electroencephalography, Euclidean alignment, test-time adaptation, batch normalisation, cross-subject, explainability.
\end{IEEEkeywords}
